\documentclass[../../main.tex]{subfiles}
\begin{document}

\begin{center}
  {\fontsize{48pt}{58pt}\selectfont \textbf{1963}\normalsize\textbf{年度}}
\end{center}

\begin{table}[htb]
  \centering
  \caption{出題テーマと難易度}
  \begin{tabular}{|c|p{9cm}|c|}
    \hline
    \textbf{問題番号} & \textbf{テーマ} & \textbf{難易度} \\
    \hline
    第1問 & 直方体の対角線と辺の長さの範囲 & \\
    \hline
    第2問 & 線分をつないだ動点の存在範囲と面積 & \\
    \hline
    第3問 & 台形内の平行線の長さと面積の無限和 & \\
    \hline
    第4問 & 正四面体に内接する三角柱の体積 & \\
    \hline
    第5問 & 面積を二等分する折線が通過する領域 & \\
    \hline
    第6問 & べき関数と放物線に挟まれた面積 & \\
    \hline
  \end{tabular}
\end{table}

\begin{table}[htb]
  \centering
  \caption{解答}
  \renewcommand{\arraystretch}{1.4}
  \begin{tabular}{|c|c|p{9cm}|}
    \hline
    \textbf{問題番号} & \textbf{小問} & \textbf{答え} \\
    \hline
    第1問 & -- & $OD=\sqrt{\dfrac{a^2+b^2+c^2}{2}}$，$4<c<\sqrt{34}$ \\
    \hline
    第2問 & -- & 範囲は解答の図を参照．面積は$4+\dfrac{2\pi}{3}$ \\
    \hline
    第3問 & (1) & $x_n=\dfrac{ab}{a+bn}$ \\
    \cline{2-3}
    & (2) & $\displaystyle\sum_{n=1}^{\infty}F_n=\dfrac{Ha^2}{2(a+b)}$（$H$は台形の高さ） \\
    \hline
    第4問 & (1) & $AP=\dfrac{2a}{3}$ \\
    \cline{2-3}
    & (2) & $\dfrac{V_0}{V}=\dfrac{4}{9}$ \\
    \hline
    第5問 & -- & $\dfrac{a^2}{3}$ \\
    \hline
    第6問 & -- & $n=3$ \\
    \hline
  \end{tabular}
\end{table}

\paragraph{試験形態}
6問．試験時間は要確認．

\end{document}
